%========================================================================
%  LIMITATIONS -- Schreibgeruest
%------------------------------------------------------------------------
%  LESEHILFE (identisch zu evaluation.tex und results.tex):
%   * "WAS HIER STEHT" = inhaltliche Checkliste, Satz fuer Satz abarbeiten.
%   * "ZAHLEN"         = die bereits VERIFIZIERTEN Werte aus dem Run.
%                        Alle Zahlen unten wurden direkt aus den Tabellen
%                        gelesen -- du musst nichts nachrechnen, aber du
%                        darfst auch KEINE Zahl schreiben, die hier nicht
%                        steht (CLAUDE.md Rigor-Regel).
%   * "ZITATE"         = biblatex-Keys. Fehlende Keys sind unten markiert.
%   * "!! ENTSCHEIDUNG"= Punkt, den DU faellen musst.
%   * "[arXiv-Kuerzung]" = wie das Unterkapitel fuer arXiv eingedampft wird.
%   * Bloecke zwischen % >>> THESIS-ONLY >>> und % <<< THESIS-ONLY <<<
%                        fliegen fuer arXiv komplett raus.
%
%  DATENHERKUNFT ALLER ZAHLEN IN DIESER DATEI
%  ------------------------------------------
%  Identisch zu results.tex: Diagnostics-Run run_full_20260810_110834
%  (Wiederholung des 03.08.-Laufs, 14 von 15 Tabellen bit-identisch --
%  Begruendung im Kopf von results.tex), Conformal-Run
%  run_full_20260721_145309, Panel-Hash 5f27a412...68ec2e,
%  git commit b1811e67. Tabellenkopien unter results/tables/*.csv.
%  AUSNAHME: die DQ-Zerlegung in 9.2 und der Punktprognose-Test in 9.3
%  sind POST-HOC-Rechnungen vom 2026-08-03, die NICHT aus den
%  Run-Artefakten stammen -- sie sind an Ort und Stelle als solche
%  markiert und muessen als Amendment deklariert werden.
%
%  GLIEDERUNGSLOGIK DIESES KAPITELS (bewusst gewaehlt, nicht willkuerlich)
%  ----------------------------------------------------------------------
%  Reihenfolge = absteigende Bedrohung fuer die Kernaussage der Arbeit:
%   9.1 Kontamination      -> kann Befund (3) VOLLSTAENDIG entwerten
%   9.2 Gate/DQ            -> relativiert das zentrale Guetesiegel
%   9.3 Inferenz/Power     -> relativiert JEDE Rangaussage
%   9.4 Seed-Fragilitaet   -> relativiert die NN-Ergebnisse speziell
%   9.5 Stichprobe/Panel   -> die harte Obergrenze von allem
%   9.6 Designentscheidungen-> was die Konstruktion praedeterminiert
%   9.7 Oekonomie          -> Proxy, kein Backtest
%   9.8 Scope              -> was die Arbeit NICHT behauptet
%  Diese Reihenfolge ist selbst ein Argument: wer zuerst die gefaehrlichste
%  Limitation nennt, wirkt nicht defensiv, sondern souveraen. NICHT
%  umsortieren, um Unangenehmes nach hinten zu schieben.
%
%  GRUNDREGEL FUER DEN TON: jede Limitation braucht (i) was genau das
%  Problem ist, (ii) die Zahl, die es belegt, (iii) in welche RICHTUNG es
%  die Ergebnisse verzerrt, (iv) was dagegen getan wurde ODER warum nichts
%  getan werden konnte. Punkt (iii) ist der, den die meisten Arbeiten
%  weglassen -- er ist der, der dich vom Gutachter unterscheidet.
%========================================================================
\section{Limitations}
\label{sec:limitations}

% WAS HIER STEHT (Kapitel-Kopf, 1 kurzer Absatz, 4-5 Saetze):
%  - Ein Satz Rahmung: die Befunde aus \cref{sec:results} gelten innerhalb
%    eines vorregistrierten Protokolls (\cref{sec:evaluation}); dieses
%    Kapitel benennt, wo das Protokoll an seine Grenzen stoesst.
%  - Ein Satz Unterscheidung, die du unbedingt machen solltest: es gibt
%    (a) Limitationen des DESIGNS (bewusste Entscheidungen mit Kosten) und
%    (b) Limitationen der EVIDENZ (was die Daten nicht hergeben). Beide
%    stehen hier, aber sie sind nicht dasselbe -- Leser verwechseln das.
%  - Ein Satz Vorwaertsverweis: die schwerwiegendste Limitation (9.1)
%    betrifft ausgerechnet das bestplatzierte Modell.
%  - KEIN Understatement ("minor caveats", "as with any study"). Solche
%    Formeln lesen sich als Verharmlosung und provozieren genau die
%    Nachfrage, die du vermeiden willst.
%  - Ein Satz, dass alle hier genannten Punkte die Rangfolge der Befunde
%    NICHT umkehren, sondern ihre Reichweite begrenzen -- sofern du das
%    nach dem Schreiben noch guten Gewissens sagen kannst. Wenn nicht:
%    weglassen, nicht abschwaechen.

% TODO Prosa Kapitelkopf.

% [arXiv-Kuerzung]: Kopf bleibt (3 Saetze).


%------------------------------------------------------------------------
\subsection{Pretraining contamination of the zero-shot models}
\label{sec:lim:contamination}
%------------------------------------------------------------------------
% WARUM DAS ZUERST KOMMT: die Kernaussage der Arbeit ist, dass ein
% Zero-Shot-Foundation-Model die besten Intervalle liefert
% (\cref{sec:res:ranking}). Genau dieses Modell ist das einzige, dessen
% Trainingsdaten du nicht kennst. Wenn diese Limitation greift, ist
% Befund (3) nicht abgeschwaecht, sondern hinfaellig. Deshalb Absatz 1.
%
% WAS HIER STEHT:
%  1. Der Sachverhalt nuechtern: TiRex-2 wurde am 2026-07-01 released; der
%     Pretraining-Cutoff ist nicht dokumentiert. Das Evaluationsfenster
%     laeuft 2015-10 bis 2026-04, endet also VOR dem Release, liegt aber
%     potenziell vollstaendig im Pretraining-Zeitraum.
%     -> Verweis \cref{sec:models:tirex2} (dort steht der Caveat kurz).
%  2. Warum das bei Staatsanleihen-Spreads besonders heikel ist: die
%     zugrundeliegenden Renditereihen sind oeffentlich, langjaehrig und
%     in nahezu jedem grossen Zeitreihen-Korpus enthalten. Es braucht
%     also keinen exotischen Datensatz, damit Kontamination plausibel ist
%     -- die Standardquellen genuegen.
%  3. WAS DAGEGEN SPRICHT (ehrlich beide Seiten, nicht nur die
%     bequeme). Drei Gegenargumente, die du machen kannst:
%     (a) Die Punktprognose ist NICHT besser als der Random Walk
%         (MAE-Ratio 0.9965, Panel-t -0.498, kein Land FDR-signifikant,
%         siehe 9.3). Ein Modell, das die Zielreihe memoriert haette,
%         muesste im Median glaenzen -- es tut es nicht.
%         -> Das ist dein STAERKSTES Gegenargument. Ausbauen.
%     (b) Der Vorsprung liegt vollstaendig in der Breite der praediktiven
%         Verteilung, nicht in ihrem Zentrum. Memorierung wuerde primaer
%         das Zentrum treffen.
%     (c) Im high-VIX-Regime verschwindet der Vorsprung (F13: tirex2 von
%         0.125 auf -0.003). Ein memorierendes Modell haette gerade in
%         den auffaelligen Stressmonaten einen Vorteil, nicht dort einen
%         Einbruch.
%  4. WAS DAFUER SPRICHT / was ungeloest bleibt: keines der drei
%     Argumente ist ein Beweis. Ohne dokumentierten Cutoff ist die Frage
%     empirisch nicht entscheidbar. Das genau so hinschreiben.
%  5. Was ein sauberer Test waere (Future Work, konkret): Evaluation auf
%     einem Fenster NACH dem Release (ab 2026-07), oder auf einer
%     Spread-Definition, die so in keiner oeffentlichen Quelle steht.
%     Beides ist hier nicht moeglich -- Fenster zu kurz bzw. ausserhalb
%     des Datenstands.
%  6. !! ENTSCHEIDUNG: wie prominent? Meine Empfehlung: Ein Satz im
%     Abstract, ein Absatz in der Introduction, dieses Unterkapitel, und
%     eine Fussnote an der Stelle in \cref{sec:res:ranking}, wo tirex2
%     zum Sieger erklaert wird. NICHT nur hier. Wenn ein Gutachter das
%     zuerst in den Limitations findet, liest es sich wie ein Eingestaendnis;
%     wenn es vorne steht, wie Souveraenitaet.
%
% ZAHLEN (verifiziert):
%   Release TiRex-2       : 2026-07-01
%   Pretraining-Cutoff    : undokumentiert
%   Evaluationsfenster    : 2015-10 bis 2026-04 (127 Monate)
%   MAE-Ratio tirex2      : 0.9965 (Median ueber Laender)
%   Panel-t Punktprognose : -0.498 (n.s.), 0/10 Laender FDR-signifikant
%   F13 tirex2            : low_vix +0.1255 -> high_vix -0.0025
%   Betroffene Modelle    : tirex, tirex2, tirex2_cov (alle ZS)
%   Nicht betroffen       : rw, arma, arma_monthly, lgbm, lstm, xlstm
%                           und deren const-Twins (lokal trainiert)
%
% ZITATE: podest2026tirex2, auer2025tirex.
%
% [arXiv-Kuerzung]: bleibt VOLLSTAENDIG. Das ist der Punkt, an dem ein
%   arXiv-Leser als Erstes haengt -- hier zu kuerzen waere ein Fehler.


%------------------------------------------------------------------------
\subsection{The admissibility gate does not test what the methods fail at}
\label{sec:lim:gate}
%------------------------------------------------------------------------
% WARUM DAS DIE WICHTIGSTE METHODISCHE LIMITATION IST:
%  Das Zulaessigkeitsgate (\cref{def:gate}) liest ausschliesslich q_UC und
%  q_CC. Beide Tests konditionieren nur auf die VERGANGENHEIT der
%  Verletzungsfolge (Kupiec auf gar nichts, Christoffersen auf I_{t-1}).
%  Der DQ-Test konditioniert zusaetzlich auf die AKTUELLE INTERVALLBREITE
%  -- und genau dieser Regressor ist es, der die Ablehnungen treibt.
%  Das Gate laesst also systematisch Kombinationen durch, deren bedingte
%  Kalibrierung entlang der Breitenachse verletzt ist.
%
% ZAHLEN (post hoc nachgerechnet 2026-08-03 auf calibrated_all.parquet des
% Runs run_full_20260721_145309, 7420 Zellen x 127 Monate = 942340 Zeilen;
% dieselbe Wald-Statistik und dieselbe BH-FDR-Familie wie im Haupt-DQ):
%   Zerlegung der DQ-Regression, Ablehnraten bei FDR q<0.10:
%     DQ voll  (const + 4 Lags + Breite, df=6)   49.3%   [= Wert aus 8.2]
%     ohne Breite (const + 4 Lags, df=5)         20.8%   (roh p<0.05: 27.7%)
%     nur Breite  (const + Breite,  df=2)        68.9%
%     Breiten-Koeffizient gamma allein, chi2(1)  53.2%   (roh p<0.05: 52.4%)
%   -> Der Satz, der hier hingehoert: NIMMT MAN DEN BREITEN-REGRESSOR
%      HERAUS, faellt die Ablehnrate von 49.3% auf 20.8%, also praktisch
%      auf das UC-Niveau (20.6%). Die DQ-Ablehnungen sind fast
%      vollstaendig ein Breiten-Phaenomen, nicht ein Lag-Phaenomen.
%      (Die Christoffersen-IND-Rate von nur 6.1% sagt dasselbe von der
%      anderen Seite: zeitliches Clustering ist NICHT das Problem.)
%
% RICHTUNG DES EFFEKTS (das ist die eigentliche Aussage):
%   Von den 3949 Zellen mit FDR-signifikantem gamma haben 99.9% ein
%   NEGATIVES gamma; Median gamma = -1.631.
%   Interpretation, die du ausschreiben musst: gamma<0 heisst, in Monaten
%   mit breitem Intervall liegt I_t - alpha unter null (Ueberdeckung), in
%   Monaten mit schmalem Intervall darueber (Unterdeckung). Die
%   Breitensteuerung reagiert in die RICHTIGE Richtung, aber sie
%   UEBERSCHIESST: sie macht zu weit auf, wenn sie aufmacht, und zu weit
%   zu, wenn sie zumacht. Marginal mittelt sich das auf ~0.80 heraus --
%   deshalb sieht Kupiec nichts. Das ist der Lehrbuchfall von bedingter
%   Fehlkalibrierung, die von marginaler Coverage maskiert wird.
%
% DER TRADE-OFF ZWISCHEN DEN METHODENKLASSEN (tab:lim_dq unten ist schon
% befuellt -- du schreibst nur Caption und Fliesstext):
%   -> Zwei saubere Fehlermodi, die du als Ergebnis verkaufen kannst:
%      (i)  STATISCHE Verfahren (native, cqr_static, saocp, decay_ocp)
%           haben geclusterte Verletzungen (Lags 40-54%) bei
%           unauffaelliger Breite (7-20%). Sie adaptieren nicht.
%      (ii) ONLINE-TRACKER (pid, spci, sfogd, pid_local, aci, agaci,
%           dtaci) beseitigen das Clustering fast vollstaendig
%           (pid_cqr: 1.5%!) und erzeugen dafuer die ueberschiessende
%           Breitenreaktion (50-97%).
%      FORMULIERUNGSVORSCHLAG: "the conformal layer trades temporal
%      dependence of violations for width-conditional dependence."
%      Das ist eine eigenstaendige Aussage ueber RQ3 und gehoert
%      zusaetzlich in \cref{sec:res:rq3} verlinkt.
%   ANMERKUNG zu mondrian_cqr: niedrigste DQ-Gesamtablehnrate aller
%      Methoden (0.291, siehe tab:res_gate) -- die Konditionierung auf
%      Volatilitaetskategorien funktioniert also genau dort, wo
%      unkonditionale Verfahren brechen. Das ist die in
%      diagnostics_foundation.txt Paragraph 6 VORREGISTRIERTE Erwartung
%      und sie trifft zu. Aber: Gate-Rate nur 0.604 und Skill -0.028 --
%      bedingt gut kalibriert, unbedingt zu breit. Ein Satz wert.
%
% DER UNBEQUEME TEIL (unbedingt hinschreiben, sonst ist es Schoenfaerberei):
%   Auf den 485 VALIDEN Kombinationen (4850 Zellen) lauten die Raten:
%     DQ voll 42.5% | gamma 65.5% | reine Lags 8.2%
%   Die gamma-Ablehnrate ist im validen Bereich also HOEHER als im
%   Gesamtdurchschnitt (65.5% vs. 53.2%). Grund: das Gate selektiert
%   bevorzugt die Online-Tracker (Gate-Rate 0.96-1.00, siehe tab:res_gate)
%   -- also genau die Methodenklasse mit dem Breitenproblem. Der Satz, der
%   hier stehen muss: "the admissibility gate does not merely fail to
%   detect this failure mode -- it actively selects for it."
%
% BETROFFENHEIT NACH MODELLFAMILIE (gamma-Ablehnrate, FDR q<0.10):
%   xlstm 71.0 | lstm 60.8 | tirex2 60.0 | tirex 51.4 | tirex2_cov 51.4 |
%   lgbm 44.3 | xlstm_const 43.8 | lstm_const 41.6 | arma_monthly 37.9 |
%   arma 36.4 | lgbm_const 35.4 | rw 30.7
%   -> Ein Satz: die Ordnung folgt der Modellkomplexitaet fast monoton,
%      und der Random Walk ist am wenigsten betroffen -- er hat schlicht
%      die geringste Breitenvariation, die falsch konditioniert sein
%      koennte. Nicht ueberinterpretieren, aber erwaehnen.
%
% !! EHRLICHKEITS-PFLICHT (sonst ist es HARKing):
%   Diese Zerlegung steht NICHT in diagnostics_foundation.txt. Sie ist
%   eine POST-HOC-Diagnostik, motiviert durch die hohe DQ-Ablehnrate im
%   Hauptlauf. Sie erzeugt kein neues Finding und aendert kein Ranking --
%   sie erklaert eine bereits vorregistriert berichtete Zahl. Trotzdem:
%   als solche deklarieren, mit Datum (2026-08-03), und in
%   diagnostics_foundation.txt Paragraph 11 als Amendment eintragen.
%   Formulierung: "post hoc decomposition of a pre-registered statistic,
%   reported to explain rather than to establish."
%
% WAS DARAUS FOLGT (Absatz 3, konstruktiv statt nur selbstkritisch):
%   1. Ein Gate, das q_DQ einschliesst, waere schaerfer -- aber es haette
%      bei n=127 kaum noch Kandidaten uebrig gelassen (42.5% der validen
%      Zellen fallen). Das ist ein echter Power-vs-Strenge-Trade-off,
%      keine Nachlaessigkeit. So framen.
%   2. Der Befund motiviert direkt zukuenftige Arbeit: breitenkonditionale
%      Rekalibrierung (Mondrian mit Breiten- statt Volatilitaetsklassen)
%      oder ein Tracker mit gedaempfter Breitenreaktion.
%   3. Fuer die Praxis: wer die Intervalle als Risikomass nutzt, bekommt
%      in ruhigen Monaten zu wenig und in bewegten zu viel Kapital
%      angesetzt -- Anschluss an den Kapital-Proxy in \cref{sec:res:econ}.
%
% ZITATE: engle2004caviar (DQ), christoffersen1998, kupiec1995,
%   gneiting2007sharpness (bedingte vs. marginale Kalibrierung),
%   vovk2003mondrian (Mondrian als Ausweg), benjamini1995fdr.
%
% [arXiv-Kuerzung]: 1 Absatz + tab:lim_dq. Die Modellfamilien-Aufschluesselung
%   und Punkt 2/3 der Folgerungen fliegen raus.

\begin{table}[t]
\centering
\caption{% TODO Caption. Decomposition of the DQ rejections: share of
         cells rejected at FDR $q<0.10$ by the width coefficient alone
         versus by the lag-only specification.}
\label{tab:lim_dq}
\begin{tabular}{lcc}
\toprule
Method & Width coefficient & Lags only \\
\midrule
\texttt{pid\_tan\_cqr}   & 0.974 & 0.325 \\
\texttt{spci\_cqr}       & 0.908 & 0.040 \\
\texttt{sfogd\_cqr}      & 0.845 & 0.170 \\
\texttt{pid\_local\_cqr} & 0.813 & 0.058 \\
\texttt{pid\_cqr}        & 0.653 & 0.015 \\
\texttt{aci\_cqr}        & 0.583 & 0.132 \\
\texttt{agaci\_cqr}      & 0.572 & 0.060 \\
\texttt{dtaci\_cqr}      & 0.521 & 0.060 \\
\texttt{cqr\_rolling}    & 0.506 & 0.126 \\
\texttt{mondrian\_cqr}   & 0.483 & 0.057 \\
\texttt{cqr\_static}     & 0.202 & 0.536 \\
\texttt{decay\_ocp\_cqr} & 0.191 & 0.404 \\
\texttt{native}          & 0.132 & 0.477 \\
\texttt{saocp\_cqr}      & 0.070 & 0.445 \\
\bottomrule
\end{tabular}
\end{table}


%------------------------------------------------------------------------
\subsection{Statistical power and the limits of the comparative inference}
\label{sec:lim:inference}
%------------------------------------------------------------------------
% WAS HIER STEHT: dieses Unterkapitel hat DREI Teile. Schreib sie als
% drei Absaetze, nicht als Liste -- sie bauen aufeinander auf.
%
% ---- ABSATZ 1: DIE RANGFOLGE IST DESKRIPTIV ----
%  Der wichtigste Satz des ganzen Kapitels, und er gehoert hierher:
%  KEINER der 66 Cross-Model-Kontraste ist FDR-signifikant. Der kleinste
%  Median-q-Wert ueber alle 66 liegt bei 0.1851.
%  ZAHLEN:
%    Staerkster Kontrast (und trotzdem n.s.):
%      XM tirex2/ZS/cqr_static vs xlstm_const/L/spci_cqr
%      d_norm -0.1156 | Panel-t -3.5488 | median-q 0.1851 | 40% der Laender
%    Zweitstaerkster:
%      XM tirex2_cov/ZS/native vs xlstm_const/L/spci_cqr
%      d_norm -0.1109 | Panel-t -4.1586 | median-q 0.1963 | 40%
%  -> Beachte die Spannung, die du benennen MUSST: der Panel-t liegt bei
%     -3.55 bzw. -4.16, also deutlich ueber jedem konventionellen
%     kritischen Wert -- und trotzdem ueberlebt der Kontrast die
%     FDR-Korrektur ueber die Laender nicht. Das ist kein Widerspruch,
%     sondern zeigt die zwei verschiedenen Fragen: der Panel-t fragt
%     "ist der Effekt im Mittel ueber die Laender da?", der Laenderanteil
%     fragt "ist er in einzelnen Laendern nachweisbar?". Erklaeren, sonst
%     wirkt es wie ein Fehler.
%  -> KONSEQUENZ, die du explizit ziehen musst: alle Rangaussagen in
%     \cref{sec:res:ranking} sind DESKRIPTIV. Formulierungen wie "X
%     outperforms Y" muessen im ganzen Dokument zu "X ranks above Y"
%     werden. Das ist eine Durchsicht-Aufgabe, keine Stilfrage.
%
% ---- ABSATZ 2: DIE MCS DISKRIMINIERT NICHT ----
%  ZAHLEN:
%    Kandidaten je Land : 26 (kuratiert, nicht alle 742)
%    Inklusionsrate     : 0.915 ueber alle Kandidaten x Laender
%    In 10/10 Laendern  : 15 von 26 Kandidaten
%    p-Werte            : Median 0.189, Maximum 0.406
%  -> Der Satz: die MCS eliminiert bei n=127 praktisch nichts. Die in
%     diagnostics_foundation.txt Paragraph 4.3 vorweggenommene
%     Nicht-Trennschaerfe ("Inklusion ~0.93") ist eingetreten, und die
%     Kuratierung des Kandidatensatzes hat sie NICHT verhindert.
%  -> Ehrlich: das ist ein Null-Befund des Verfahrens, nicht der Modelle.
%     Er sagt nicht "alle Modelle sind gleich gut", sondern "diese
%     Stichprobe kann sie nicht trennen". Der Unterschied ist wichtig.
%
% ---- ABSATZ 3: DM-VARIANZ MIT HAC-LAG 0 ----
%  WAS DER PUNKT IST:
%   Der paarweise DM-Test (\cref{sec:eval:comparison}) schaetzt die
%   Varianz des mittleren Verlustdifferentials mit der einfachen
%   Stichprobenvarianz, also HAC-Lag 0. Die Begruendung ist
%   lehrbuchkonform: bei Prognosehorizont h sind die Fehlerdifferentiale
%   mechanisch bis Lag h-1 autokorreliert, und bei h=1 ist h-1=0.
%   In diesem Design ist der Horizont aber NICHT die einzige Quelle
%   serieller Korrelation:
%     (i)  Refit alle 12 Monate: alle Prognosen innerhalb eines Blocks
%          teilen denselben Parametervektor, ihre Fehler sind daher
%          innerhalb des Blocks korreliert (siehe tab:refit_blocks).
%     (ii) Die adaptiven CP-Methoden tragen Zustand ueber die Zeit
%          (ACI-Fehlerintegral, PID-Integralterm, SAOCP-Expertengewichte);
%          zwei aufeinanderfolgende Intervalle sind daher nicht
%          unabhaengig konstruiert.
%   Beide Kanaele erzeugen positive Autokorrelation in d_t, die die
%   Lag-0-Varianz ignoriert. Unterschaetzte Varianz -> aufgeblaehte
%   |t|-Werte -> ZU KLEINE p-Werte. Die Einzel-Land-DM-p-Werte sind also
%   systematisch zu optimistisch, und zwar in bekannter Richtung.
%
%  WAS DAGEGEN TATSAECHLICH GETAN WURDE (nicht nur vorgeschlagen --
%  implementiert, Code: diagnostics_master_v2.ipynb, Zelle 15, Funktionen
%  _nw_tstat und run_contrast):
%   Der zeitgeclusterte Panel-t (\cref{sec:eval:comparison}) ist die
%   Absicherung. Zweistufig:
%     Schritt 1: die pro Land skalenfrei normierten Differentiale
%       d_{i,t}/Wbar_i^ref werden JE MONAT ueber die 10 Laender gemittelt.
%       Das buendelt den gemeinsamen Schock in EINE Monatszahl, statt ihn
%       zehnmal zu zaehlen (Querschnittskorrelation der Euro-Spreads).
%     Schritt 2: auf dieser 127-Monats-Reihe wird der t-Wert des
%       Mittelwerts mit Newey-West-Langfristvarianz gebildet
%       (Bartlett-Kernel, Lag 3):
%       sigma^2_LR = gamma_0 + 2*sum_{l=1}^{3} (1 - l/4) * gamma_l.
%       Das faengt genau die Rest-Autokorrelation ab, die Lag 0 verpasst.
%   Das ist die praktikable Ein-Achsen-Variante der Driscoll-Kraay-Logik
%   (Querschnitt-Mittelung + HAC ueber Zeit). WICHTIG fuer die
%   Ehrlichkeit: sie beansprucht KEINE volle zweidimensionale
%   Driscoll-Kraay-Kovarianz -- das explizit hinschreiben.
%   -> Spalte t_panel in dm_contrasts.csv; auf ihr beruhen alle
%      RQ-Aussagen in \cref{sec:res:rq3} bis \cref{sec:res:rq2}.
%
%  WAS TROTZDEM OFFEN BLEIBT (das ist die eigentliche Limitation):
%   1. Die PRO-LAND-DM-Tests selbst bleiben unkorrigiert. Aus ihnen
%      entstehen die q-Werte, und aus den q-Werten die Spalte
%      share_sig_fdr ("Anteil FDR-signifikanter Laender"), die in jeder
%      Kontrast-Tabelle steht. Diese Anteile sind daher nach oben
%      verzerrt. Wo du sie berichtest, gehoert ein Halbsatz dazu.
%   2. Lag 3 ist eine Konvention, kein datengetriebener Wert (kein
%      Newey-West-Andrews-Plugin, keine Sensitivitaetsanalyse ueber
%      Lag 0/3/6/12). Als Wahl deklarieren, nicht als Optimum.
%   3. Der Panel-t setzt implizit voraus, dass die Querschnittsmittelung
%      die Laenderkorrelation VOLLSTAENDIG absorbiert. Bei stark
%      heterogenen Laendern (mittleres Kendall-tau 0.346, siehe 9.5)
%      ist das eine Naeherung.
%
%  PRAKTISCHE KONSEQUENZ -- und warum sie hier klein ausfaellt:
%   Der wichtigste Cross-Model-Befund ist ohnehin ein NICHT-Befund
%   (Absatz 1). Eine Verzerrung, die p-Werte ZU KLEIN macht, kann diesen
%   Befund nicht erzeugt haben -- sie wuerde in die Gegenrichtung wirken.
%   Der Nicht-Befund ist also robust gegen genau diese Limitation. Das
%   ist ein starkes Argument und gehoert hingeschrieben. Betroffen sind
%   dagegen die RQ1/RQ2-Kontraste, wo Signifikanz BEHAUPTET wird (z.B.
%   RQ2 lstm/GF native: t_panel 8.18, share_sig_fdr 0.633) -- dort ist
%   der Panel-t die tragende Zahl, nicht der Laenderanteil.
%
% >>> THESIS-ONLY >>>
% - Absatz zur Power-Rechnung: bei n=127 und alpha=0.20 erwartet man
%   25.4 Verletzungen. Um eine Coverage von 0.75 gegen 0.80 auf 5% Niveau
%   zu trennen, braucht der Kupiec-LR-Test deutlich mehr Beobachtungen.
%   Wenn du das rechnest, rechne es SAUBER und schreib die Annahmen dazu
%   -- sonst weglassen. Es ist kein Pflichtteil.
% - Absatz zum Verhaeltnis Panel-t vs. Laenderanteil an einem konkreten
%   Beispiel durchgerechnet.
% <<< THESIS-ONLY <<<
%
% ZITATE: dieboldmariano1995, harvey1997, newey1987, hansen2011mcs,
%   kunsch1989mbb, benjamini1995fdr, giacomini2006 (dokumentierte
%   Scope-Entscheidung: bei h=1, festem Refit-Protokoll und identischen
%   Informationsmengen faellt der GW-Vorteil mit diesem Design weitgehend
%   zusammen -- bewusst nicht implementiert).
%   !! FEHLEN NOCH in references.bib: driscollkraay1998, petersen2009.
%      Beide werden fuer den Panel-t-Absatz gebraucht -- eintragen.
%
% [arXiv-Kuerzung]: Absatz 1 bleibt vollstaendig (er ist die
%   Kernrelativierung), Absatz 2 auf 2 Saetze, Absatz 3 in eine Fussnote
%   zum DM-Absatz in \cref{sec:eval:comparison}.


%------------------------------------------------------------------------
\subsection{Seed fragility of the neural results}
\label{sec:lim:seeds}
%------------------------------------------------------------------------
% WAS HIER STEHT:
%  1. Der Sachverhalt: fuer xLSTM und LSTM laeuft die komplette Pipeline
%     pro Seed (42/43/44); Prognosen werden nie ueber Seeds gemischt
%     (\cref{sec:exp:seeds}). Die Seed-Spanne des Skills ist damit ein
%     direktes Mass fuer den Zufallsanteil im Ergebnis.
%  2. DIE ZAHL, die weh tut: bei 51.7% der 588 NN-Kombinationen ist die
%     Seed-Spanne des Median-Skills GROESSER als der Betrag des
%     Median-Skills selbst. Bei ueber der Haelfte der neuronalen
%     Ergebnisse ist die Initialisierung also einflussreicher als der
%     gemessene Effekt.
%  3. Der Fall, den du namentlich nennen musst, weil er das beste
%     xLSTM-Paket ist: xlstm/L/native hat Median-Skill 0.0939 bei einer
%     Seed-Spanne von 0.1022. Das beste Ergebnis der namensgebenden
%     Architektur liegt innerhalb seiner eigenen Seed-Streuung.
%     -> Das nicht relativieren. Es ist ein sauber gemessenes negatives
%        Resultat und gehoert genau so berichtet.
%  4. Was das Protokoll dagegen tut: die Ex-ante-Regel aus
%     diagnostics_foundation.txt Paragraph 4.4 verbietet eine Rangaussage,
%     wenn das Vorzeichen des Kontrasts nicht ueber alle Seeds stabil ist;
%     betroffene Faelle werden als F9 ausgewiesen und die Aussage
%     unterbleibt. Das ist die richtige Reaktion, aber sie heilt das
%     Problem nicht, sie macht es nur sichtbar.
%  5. Was NICHT getan wurde und warum: drei Seeds sind wenig. Fuenf oder
%     zehn waeren belastbarer, kosten aber die komplette Kette
%     (Modell -> CP -> Diagnostik) pro zusaetzlichem Seed. Als bewusste
%     Ressourcenentscheidung deklarieren, nicht als Versehen.
%  6. Ein Satz zur Richtung der Verzerrung: Seed-Rauschen ist
%     unsystematisch, es verzerrt den Median ueber Kombinationen also
%     nicht in eine Richtung -- aber es blaeht die Streuung auf und macht
%     einzelne Spitzenwerte unzuverlaessig. Wer den BESTEN Seed
%     berichtete, wuerde nach oben verzerren; hier wird nicht selektiert.
%
% ZAHLEN (verifiziert, ranking_all.csv, Spalte seed_span):
%   NN-Kombinationen gesamt (lstm/xlstm inkl. const-Twins) : 588
%   davon seed_span > |median_skill|                       : 51.7%
%   xlstm/L/native : median_skill 0.0939, seed_span 0.1022
%   Seeds          : 42, 43, 44 (nur NN; lgbm deterministisch mit
%                    deterministic=True, rw/arma/tirex/tirex2 ohnehin)
%   F9-Findings im Story-Report: 3 Instanzen, Score je 0.44
%
% ZITATE: keine Pflichtzitate. Optional eine Quelle zur
%   Seed-Varianz in Deep-Learning-Benchmarks, falls du eine hast --
%   sonst weglassen, die eigene Zahl traegt den Punkt.
%
% [arXiv-Kuerzung]: 2 Saetze (die 51.7% und der xlstm/L/native-Fall).
%   Diese beiden Zahlen MUESSEN drin bleiben.


%------------------------------------------------------------------------
\subsection{Sample size, cross-section and country heterogeneity}
\label{sec:lim:sample}
%------------------------------------------------------------------------
% WAS HIER STEHT: das ist die "harte Obergrenze"-Sektion. Vier Punkte.
%
%  1. STICHPROBE. Panel 10 Laender x 277 Monate = 2770 Beobachtungen;
%     davon 127 Monate je Land im Test. Das ist die Obergrenze fuer
%     alles, was hier gemessen werden kann, und der gemeinsame Nenner
%     von 9.3 (Power) und 9.4 (Seeds).
%     -> Ein Satz, der die Verbindung zum xLSTM-Ergebnis zieht:
%        277 Monate sind fuer eine Architektur dieser Kapazitaet wenig.
%        Das ist die plausibelste Erklaerung fuer das negative
%        Kernergebnis und gehoert hier statt in \cref{sec:res:ranking}.
%
%  2. QUERSCHNITTSKORRELATION. Zehn Euro-Spreads sind keine zehn
%     unabhaengigen Evidenzstuecke. Adressiert per Panel-t (9.3), aber
%     nicht eliminiert. Die effektive Zahl unabhaengiger Beobachtungen
%     ist unbekannt und kleiner als 1270.
%
%  3. LAENDERHETEROGENITAET. Das mittlere paarweise Kendall-tau zwischen
%     den Kombinations-Rankings der Laender betraegt 0.346.
%     ZAHLEN (verifiziert, aus cells.csv, pool=rolling, Pivot
%     Kombination x Land auf skill; 45 Laenderpaare):
%       Mittel 0.3457 | Median 0.3436 | Min 0.0734 | Max 0.6496
%       Niedrigste Paare : FI-GR 0.073 | IE-IT 0.082 | GR-IT 0.097
%       Hoechste Paare   : ES-PT 0.650 | GR-PT 0.623 | ES-GR 0.563
%       Mittleres tau je Land (am wenigsten anschlussfaehig zuerst):
%         IT 0.225 | FI 0.280 | FR 0.335 | NL 0.341 | IE 0.341 |
%         GR 0.345 | AT 0.373 | PT 0.392 | ES 0.408 | BE 0.418
%     -> DIE INTERPRETATION, die du ziehen musst: die hoechsten
%        Korrelationen liegen zwischen den Peripherielaendern
%        (ES/PT/GR), die niedrigsten zwischen Kern und Peripherie
%        (FI-GR 0.073). Es gibt also nicht EINE Rangordnung, sondern
%        mindestens zwei Regime im Querschnitt. Medianaussagen ueber
%        alle zehn Laender bleiben gueltig, aber "das beste Verfahren"
%        ist laenderabhaengig. Das relativiert jede Einzelempfehlung.
%     -> QUERVERBINDUNG, die gratis ist: die Tail-Asymmetrie streut
%        ueber Laender staerker (0.431 bis 0.642) als ueber Modelle
%        (0.483 bis 0.618), siehe \cref{sec:res:ranking}. Zwei
%        unabhaengige Kennzahlen sagen dasselbe. Explizit verknuepfen.
%
%  4. EINZELNE LAENDER. Zwei Faelle, die eigene Saetze brauchen:
%     (a) GRIECHENLAND. Gate-Rate nur 0.519 -- das mit Abstand am
%         schwersten kalibrierbare Land. Dazu der PSI-Strukturbruch 2012
%         (vgl. \cref{sec:data:universe}) und die fehlenden Monate
%         2015-07/08 (Kapitalkontrollen, liegen im Burn-in, nicht im
%         Test). WICHTIG: es gibt EINEN Lauf inklusive GR und KEINEN
%         ex-GR-Robustheitscheck. Das ehrlich als offene Flanke
%         benennen und als Future Work markieren -- nicht so tun, als
%         waere es geprueft.
%     (b) IRLAND. Der bewertungsseitige ex-IE-Check
%         (\cref{sec:eval:robustness}) zeigt: die ORDNUNG des Rankings
%         ist stabil (Spearman-rho 0.9630, Kendall-tau 0.8493), das
%         NIVEAU nicht. Der Median-Skill des Siegers faellt von 0.1461
%         auf 0.1027, wenn Irland aus der Aggregation genommen wird --
%         ein Drittel des Vorsprungs haengt an einem Land.
%         Median d_skill ueber alle Kombinationen -0.0118; bei 35.6%
%         der Kombinationen aendert sich der Skill um mehr als 0.02.
%         -> Der Satz: Rangaussagen sind robust, Effektgroessen nicht.
%            Wo du eine Skill-ZAHL nennst, gehoert diese Sensitivitaet
%            dazu.
%     Gate-Rate je Land vollstaendig (calibration_tests.csv, cell_valid):
%       GR 0.519 | PT 0.616 | IE 0.623 | ES 0.660 | BE 0.718 |
%       AT 0.747 | NL 0.771 | IT 0.810 | FI 0.844 | FR 0.864
%       -> Ein Satz Interpretation: die Peripherie mit sprunghaften
%          Spreads ist am schwersten kalibrierbar, und zwar monoton.
%
% ZITATE: kendall1938 (tau), benjamini1995fdr.
%   !! kendall1938 FEHLT in references.bib -- eintragen.
%      driscollkraay1998 und petersen2009 ebenfalls (siehe 9.3).
%
% [arXiv-Kuerzung]: Punkte 1-3 auf einen Absatz, Punkt 4 auf zwei Saetze
%   (GR-Gate-Rate und der ex-IE-Niveaueffekt 0.146 -> 0.103). Der
%   ex-IE-Effekt MUSS drin bleiben.


%------------------------------------------------------------------------
\subsection{Design decisions that bound the results}
\label{sec:lim:design}
%------------------------------------------------------------------------
% WAS HIER STEHT: Limitationen, die aus BEWUSSTEN Entscheidungen folgen.
% Abgrenzung zum Rest des Kapitels im ersten Satz klarmachen: hier geht
% es nicht um fehlende Evidenz, sondern um Kosten getroffener Wahlen.
% Sechs Punkte, je 2-3 Saetze. Reihenfolge nach Tragweite.
%
%  1. SYMMETRISCHER CQR-SCORE. Der Score E = max(q_lo - y, y - q_hi) ist
%     einseitig definiert und behandelt beide Intervallraender gleich.
%     Die Daten sind aber nicht symmetrisch: bei praktisch gleicher
%     ANZAHL von Verletzungen oben und unten (Tail-Asymmetrie im Mittel
%     0.5049) ist die Ueberschreitungsstrafe deutlich groesser als die
%     Unterschreitungsstrafe (0.0715 gegen 0.0521, also +37%).
%     -> Wenn der Spread aus dem Intervall laeuft, dann nach oben und
%        weiter. Ein symmetrischer Score kann das strukturell nicht
%        abbilden. Als Future Work benennen (asymmetrische oder
%        zweiseitig getrennte Score-Familien).
%     -> Nicht kausal formulieren, nur als Beobachtung, die zur
%        bekannten Rechtsschiefe von Kreditspreads passt.
%
%  2. BURN-IN-DOPPELROLLE. Das Validierungsfenster dient sowohl der
%     Hyperparameterwahl als auch der Initialisierung des
%     Kalibrierpools. Diagnostiziert (F11), Effekt klein und
%     selbstlimitierend:
%       Coverage erste 36 Testmonate minus Rest, Mittel ueber alle
%       Kombinationen: -1.07pp; groesste Einzeldifferenz bei native
%       mit -6.64pp.
%     -> Ein Satz: der Anfangsoptimismus existiert, ist auf die nicht
%        rekalibrierten Streams konzentriert und verschwindet mit
%        wachsendem Pool.
%
%  3. KALIBRIERPOOL-FENSTER. Primaerfall rolling 36; die Varianten
%     rolling24/rolling48/expanding sind Robustheitspaesse und gelten
%     nur fuer pool-lesende Methoden (Online-Tracker lesen den Pool
%     nicht -- deren Spalten sind in den Varianten identisch).
%     ZAHLEN (cells.csv, mittlere Coverage / Median-Skill je pool_mode):
%       expanding 0.8432 / -0.0111 | rolling  0.8341 / -0.0126 |
%       rolling24 0.8127 /  0.0003 | rolling48 0.8204 /  0.0142
%     -> Die Fenstergroesse verschiebt die Coverage um bis zu 3pp. Das
%        ist kein Nebeneffekt, sondern eine echte Designabhaengigkeit.
%     -> F8-Bruch mit dem groessten Ausschlag: lstm/GI/cqr_rolling mit
%        8.9pp Coverage-Differenz rolling vs. expanding.
%     -> EHRLICH: die vorregistrierte Erwartung war, dass SPCI wegen
%        Pool-Duennbesetzung bricht. Getroffen hat es cqr_rolling.
%        SPCI ist mit 0.769 (rolling) vs 0.795 (expanding)
%        vergleichsweise stabil. Die Vorhersage war FALSCH und das
%        gehoert so berichtet -- eine widerlegte Vorregistrierung ist
%        ein Qualitaetsmerkmal, kein Makel.
%
%  4. QUANTILKREUZUNG. Im Mittel vernachlaessigbar (0.451% aller
%     kalibrierten Intervalle; 81.7% der Zellen kreuzen nie; auf den
%     validen Kombinationen 0.398%), aber ungleich verteilt:
%       pid_tan_cqr 2.554% | cqr_static 0.724% mit Extremfaellen
%       bis 77.95% (xlstm/L/seed42/IE); 18 von 530 cqr_static-Zellen
%       liegen ueber 5%.
%     -> Mechanismus in \cref{sec:conformal:notes} erklaeren, hier nur
%        die Konsequenz: cqr_static ist die Methode mit dem hoechsten
%        Rohskill, und tirex2/ZS/cqr_static hat in PT eine Kreuzungsrate
%        von 32.3%. Das ist seit 2026-08-10 ein ARGUMENT FUER die in
%        \cref{sec:res:ranking} gefallene Entscheidung (Headline-Paket ist
%        decay_ocp_cqr mit 0.028%), nicht mehr ein Caveat GEGEN den
%        Sieger -- so formulieren, sonst widerspricht der Absatz 8.3.
%
%  5. REGIME-KONVENTIONEN. Zwei Festlegungen, die Ergebnisse mitbestimmen:
%     (a) Die ersten 12 Testmonate haben keinen Volatilitaetswert (das
%         rollierende Fenster sieht nur Testrealisationen) und werden
%         neutral als "mid" gefuehrt. Konvention, dokumentiert per
%         Amendment 2026-07-10.
%     (b) Der VIX-Schwellenwert 20 ist ex ante fixiert, aber
%         konventionell gewaehlt; VIX>25 wird nur nachrichtlich
%         mitgefuehrt.
%     -> Ein Satz: beide Festlegungen sind vor der Auswertung getroffen,
%        aber keine ist datengetrieben optimiert. Genau so schreiben.
%     -> WICHTIG UND UNBEQUEM: die beiden Stressdefinitionen liefern
%        GEGENSAETZLICHE Vorzeichen. Ueber die Volatilitaetsterzile
%        STEIGT die Coverage im Stress (calm 0.7795 -> stressed 0.8412),
%        ueber den VIX FAELLT sie (low 0.8342 -> high 0.7588). Das sind
%        dieselben Daten unter zwei Stratifizierungen. Erklaerbar (das
%        Vol-Terzil ist laenderspezifisch und rueckwaertsgewandt, der
%        VIX global und gleichzeitig), aber du MUSST es adressieren --
%        sonst fragt es der Gutachter.
%
%  6. FESTE SPLITS UND REFIT-FREQUENZ. Der 12-Monats-Refit ist empirisch
%     abgesichert (F10: Median-Effekt ueber 14 Kontraste 0.65%, nur 14%
%     FDR-signifikant), aber nur GEGEN monatliches Refit und nur fuer
%     ARMA. Fuer die NN-Modelle gibt es keinen entsprechenden Check --
%     der Rechenaufwand waere prohibitiv. Als Luecke benennen.
%     -> Ausnahme, die du nennen solltest: bei cqr_static ist der
%        Refit-Effekt mit 6.4% (Panel-t 39.15) deutlich. Der statische
%        Kalibrierpool reagiert auf die Refit-Frequenz, die adaptiven
%        Verfahren nicht. Konsistent und erklaerbar.
%
% ZITATE: romano2019cqr (Symmetrie des CQR-Scores), winkler1972,
%   gneiting2007, vovk2003mondrian, xu2023spci.
%
% [arXiv-Kuerzung]: Punkte 1, 3 und 5(b) bleiben; 2, 4, 6 in einen
%   Sammelsatz. Der Widerspruch der beiden Stressdefinitionen MUSS
%   drin bleiben.


%------------------------------------------------------------------------
\subsection{The economic evaluation is a proxy, not a backtest}
\label{sec:lim:econ}
%------------------------------------------------------------------------
% WAS HIER STEHT: kurz halten (1 Absatz), aber praezise. Der Abschnitt
% \cref{sec:eval:economic} rahmt die Oekonomie bereits als ergaenzende
% Evidenz -- hier steht, was das konkret ausschliesst.
%
%  1. Die Handelsregel hat KEINE Transaktionskosten, keine Bid-Ask-
%     Spannen, keine Finanzierungskosten und keine Kapazitaetsgrenzen.
%     Sie ist ausschliesslich KOMPARATIV zu lesen, nie als erzielbare
%     Rendite.
%  2. Die tatsaechliche Aktivitaet ist minimal: der Signalfilter
%     (Position nur wenn 0 nicht im Intervall liegt) laesst im Median
%     nur rund 2% der Monate zu Trades werden. Ein Sharpe-Wert, der auf
%     so wenigen aktiven Monaten beruht, ist statistisch praktisch
%     bedeutungslos.
%     ZAHLEN: Median-Sharpe ueber valide Kombinationen 0.159; 67.6%
%     positiv; Maximum 0.621 (lstm/L/pid_cqr) bei active_share 0.020.
%  3. Der Kapitalproxy setzt "Intervallbreite proportional zur
%     Kapitalunterlegung" -- eine dokumentierte Vereinfachung der
%     FRTB-Logik, keine Umsetzung der Regeltexte.
%  4. DER INTERESSANTE TEIL, den du nicht weglassen solltest, weil er
%     die beiden Kennzahlen gegeneinander stellt:
%       corr(Median-Skill, Kapitalquotient) = -0.711
%       corr(Median-Skill, Median-Sharpe)   = -0.147
%     -> Intervallguete zahlt auf die Kapitalkennzahl ein, aber kaum auf
%        die Handelsperformance. Anders gesagt: der Nutzen besserer
%        Intervalle liegt im Risikomanagement, nicht im Alpha. Das ist
%        eine substanzielle Aussage und rechtfertigt, dass der
%        Oekonomie-Block ueberhaupt existiert.
%     -> Vorzeichen erklaeren: ein NIEDRIGER Kapitalquotient ist gut
%        (engere valide Intervalle), deshalb ist die negative
%        Korrelation mit dem Skill der erwuenschte Zusammenhang.
%        Unbedingt dazuschreiben, sonst liest es sich falschherum.
%
% ZAHLEN AUS: results/tables/econ_results.csv (742 x 7), Spalten
%   median_sharpe, active_share, median_cap_ratio; gejoint auf
%   ranking_all.csv ueber (model, regime, seed, cp_method), gefiltert
%   auf validity_flag.
%
% ZITATE: bcbs2019mar (FRTB).
%   !! FEHLT in references.bib: leitchtanner1991 (Leitch & Tanner 1991,
%      "Economic forecast evaluation: profits versus the conventional
%      error measures", AER 81(3), 580-590). Wird in
%      diagnostics_foundation.txt Paragraph 7 als Begruendung des ganzen
%      Blocks zitiert -- gehoert hier her.
%
% [arXiv-Kuerzung]: 3 Saetze. Die beiden Korrelationen bleiben, der Rest
%   in einen Halbsatz.


%------------------------------------------------------------------------
\subsection{Scope: what this study does not claim}
\label{sec:lim:scope}
%------------------------------------------------------------------------
% WAS HIER STEHT: der Abschluss. Kein Jammern, sondern eine praezise
% Abgrenzung. Fuenf Punkte, je 1-2 Saetze. Das ist die Sektion, die dem
% Gutachter zeigt, dass du weisst, was du gemessen hast.
%
%  1. KEINE KAUSALITAET. Alle Befunde sind praediktiv-komparativ
%     innerhalb dieses Protokolls. Kein Feature-Effekt wird kausal
%     interpretiert, auch nicht der Kovariaten-Kontrast
%     tirex2_cov - tirex2.
%
%  2. KEINE PUNKTPROGNOSE-UEBERLEGENHEIT. Das Framework testet
%     MAE-vs-RW nicht als vorregistrierte Hypothese; die MAE-Ratios sind
%     deskriptiv (\cref{sec:eval:skill}, Punktdimension).
%     -> POST-HOC NACHGETRAGEN 2026-08-03 und als Amendment zu
%        deklarieren: DM/HLN auf |e|-Differentialen plus Panel-t
%        (NW Lag 3) fuer tirex2/ZS gegen rw/L ergibt Panel-t -0.498,
%        0 von 10 Laendern FDR-signifikant, 5 von 10 Laendern mit
%        MAE-Ratio unter 1. Der Vorsprung von 0.35% ist Rauschen.
%     -> KONTROLLE, die den Test wertvoll macht: dieselbe Prozedur auf
%        arma/L gegen rw/L ergibt Panel-t +2.447 (arma signifikant
%        SCHLECHTER, FI mit q 0.0040). Der Test hat also Power, wenn ein
%        Effekt vorhanden ist. Bei tirex2 findet er keinen, weil keiner
%        da ist. Das ist eine echte Nicht-Ablehnung, keine
%        Power-Schwaeche -- diesen Kontrollbefund unbedingt mitschreiben.
%     -> WEITERER POST-HOC-BEFUND, der eine offene Frage aus
%        \cref{sec:res:point} schliesst: die Richtungstrefferquote des
%        Random Walk (0.5827, hoechster Wert im Feld) ist ein ARTEFAKT.
%        Die RW-Delta-Prognose ist konstant null, damit wird der Nenner
%        der PT-Statistik exakt null und der Test ist undefiniert --
%        pt_stat ist in 90% der rw-Zellen NaN. Die Zahl ist schlicht der
%        Anteil der Monate ohne Spread-Anstieg, und "0% signifikante
%        PT-Tests" heisst "Test undefiniert", nicht "keine Evidenz".
%        Als Artefakt benennen, nicht deuten.
%        Zum Vergleich das Gesamtbild: mittlere Trefferquote 0.5397,
%        7.74% der DEFINIERTEN PT-Tests auf 5% signifikant (bei reinem
%        Zufall erwartet man 5%). Praktisch keine ausbeutbare
%        Richtungsinformation.
%
%  3. F12 IST LEER -- und zwar per Konstruktion. Der vorregistrierte
%     Finding-Typ "Punkt- vs. Intervalldivergenz" verlangt
%     MAE-Ratio < 0.98 bei negativem Skill. Das beste MAE-Verhaeltnis im
%     gesamten Feld ist 0.9965. Die Bedingung kann nicht erfuellt werden.
%     -> Das AKTIV berichten, nicht verschweigen. Ein leerer
%        vorregistrierter Finding-Typ ist ein Beleg dafuer, dass der
%        Katalog vorher feststand und nicht nachtraeglich passend
%        gemacht wurde. Von den 13 Typen F1-F13 haben 12 Instanzen
%        erzeugt, F12 keine.
%
%  4. GIACOMINI-WHITE NICHT IMPLEMENTIERT. Dokumentierte
%     Scope-Entscheidung, kein Versehen: bei h=1, festem Refit-Protokoll
%     und identischen Informationsmengen faellt der GW-Vorteil
%     (bedingte Prognosefaehigkeit, Schaetzunsicherheit) mit diesem
%     Design weitgehend zusammen. Begruendung in
%     diagnostics_foundation.txt Paragraph 4.1.
%
%  5. NACHTRAEGLICHE ERGAENZUNGEN, vollstaendig deklariert. Fuer die
%     Nachvollziehbarkeit alle Post-hoc-Rechnungen an EINER Stelle
%     auflisten, mit Datum und der Aussage, ob sie ein Finding erzeugen:
%       (a) DQ-Zerlegung (9.2), 2026-08-03 -- erklaert eine
%           vorregistrierte Zahl, erzeugt kein Finding, aendert kein
%           Ranking.
%       (b) Punktprognose-DM gegen RW (9.3 / Punkt 2 oben), 2026-08-03 --
%           erzeugt kein Finding, bestaetigt einen deskriptiv bereits
%           berichteten Nicht-Befund.
%       (c) Laender-Kendall-tau (9.5) -- war als F7 vorregistriert, aber
%           in \cref{sec:eval:robustness} nicht eingefuehrt; die Methode
%           ist nachgetragen, die Zahl stammt aus dem Hauptlauf.
%     -> Alle drei gehoeren auch in diagnostics_foundation.txt
%        Paragraph 11 (Amendments) mit demselben Datum. Der Text hier
%        und das Protokoll dort muessen uebereinstimmen.
%
% >>> THESIS-ONLY >>>
% - Ein Absatz zur Reproduzierbarkeit: identischer data_sha256 ueber
%   alle sieben Modellfamilien, git-Commit, Versionsstand, Laufzeiten.
%   Das gehoert eigentlich in den Anhang, passt aber als Gegengewicht
%   ans Ende der Limitations -- es zeigt, was SEHR wohl abgesichert ist.
% <<< THESIS-ONLY <<<
%
% ZITATE: giacomini2006, pesaran1992, benjamini1995fdr.
%
% [arXiv-Kuerzung]: Punkte 1-3 auf einen Absatz, 4-5 in eine Fussnote.
%   Punkt 2 (kein Punktprognose-Vorsprung) MUSS drin bleiben -- er ist
%   das Gegengewicht zur Kontaminationsdebatte in 9.1.
